% Introduction — Securities Lending and Information Acquisition

When a mutual fund lends shares from its portfolio to a short seller, does the fund learn something valuable?  Short sellers are among the best-informed participants in financial markets, and a large body of evidence shows that their trades predict negative future returns.\footnote{See, among others, \cite{diamond_verrecchia:87}, \cite{boehmer_jones_zhang:08}, \cite{engelberg_reed_ringgenberg:2012}, and \cite{karpoff_lou:10}.}  By participating in the securities lending market, fund managers observe, in real time and at the security level, which of their holdings are being borrowed.  This information is private: aggregate short interest is published only with a delay, and position-level lending data are not publicly available.  We ask whether active fund managers exploit this private information channel by adjusting their portfolios in response to lending activity, and, if so, through what mechanism the information is transmitted and acted upon.

We study this question using a unique regulatory dataset from the Deutsche Bundesbank that records the security-level lending positions of all German open-end equity mutual funds.  Our sample covers 320 funds from 36 fund families over the period December 2014 to December 2018, during which the European Union's short position disclosure regime (\citealp{jones_reed_waller:16}) provides an additional source of identification.  For each fund-security-month observation, we observe both the fraction of the position on loan and the change in the number of shares held, allowing us to trace how lending activity maps into subsequent portfolio adjustments at a granularity that has not previously been available in regulatory data.

Our main empirical findings are as follows.

\emph{First}, active funds reduce their positions in stocks that they lend.  A one-unit increase in the share of a position on loan is associated with a 0.33 percentage point reduction in the position over the following month (Table~\ref{table_baseline}, column~3; $t$-statistic $>$ 3).  Given a median position weight of 0.63\%, this effect is economically large, amounting to more than half of the typical holding.

\emph{Second}, the effect is concentrated among active funds.  When we interact the on-loan indicator with a dummy for active management, the interaction coefficient is $-0.34$ percentage points and highly significant, while the baseline effect for passive funds is indistinguishable from zero (Table~\ref{table_baseline}, column~2).  Passive funds, which are constrained to track an index and cannot freely adjust their holdings, serve as a natural placebo group.

\emph{Third}, we confirm the information acquisition channel by exploiting the European Union's mandatory short position disclosure regime.  Under this regime, any investor holding a net short position exceeding 0.5\% of shares outstanding must publicly disclose the position by the next business day.  If lending conveys private information about impending short-sale demand, then the informativeness of the lending signal should be highest when a public disclosure event---validating the signal---follows shortly thereafter.  We find precisely this pattern: the interaction between lending and a lead disclosure indicator is $-0.80$ to $-0.86$ percentage points ($p < 0.05$), implying that the effect of lending on position changes more than doubles when a mandatory disclosure event occurs in the subsequent month (Table~\ref{table_disclosure}).  Passive funds again show no response.

\emph{Fourth}, the timing of the disclosure event matters in the predicted direction.  Only the \emph{lead} disclosure---reflecting future information arrival---amplifies the lending effect.  Contemporaneous and lagged disclosure interactions are statistically insignificant (Table~\ref{table_timing}), ruling out a mechanical explanation and confirming that active managers act on information embedded in current lending before it becomes publicly available.

\emph{Fifth}, we address the concern that lending and trading decisions are jointly determined by using passive fund lending as an exogenous source of variation in lending activity.  Because passive funds do not adjust their portfolios in response to lending signals, their lending behavior is driven by supply-side factors rather than information.  Our results are robust to this specification (Table~\ref{table_passive_iv}).

\emph{Sixth}, information flows are strongest within fund families.  Active funds reduce their positions more aggressively when another fund managed by the same portfolio manager lends the same security, consistent with information sharing through intra-family channels (Table~\ref{table_family}).  This finding is consistent with the fund family organization literature (\citealp{evans_porrasprado_zambrana:2020}) and suggests that lending desks serve as conduits of private information within management companies.

\emph{Seventh}, the effect is amplified in concentrated lending markets---those with fewer active agents on the supply side (Table~\ref{table_market_structure}).  In such markets, each lending transaction carries more information, as the signal-to-noise ratio is higher.

These findings speak to a central tension in the securities lending literature.  \cite{EvansFerreiraPrado2017} document that active funds that participate in lending underperform otherwise similar non-lending funds by 0.5\%--0.7\% per year, a puzzle they attribute to fund family incentives that prevent managers from selling stocks with high short-selling demand.  Our results offer a partial resolution: active managers \emph{do} respond to the lending signal by reducing positions, but the response is concentrated among managers who face fewer institutional constraints, such as those operating in active mandates with discretion over portfolio weights.  The underperformance documented by \cite{EvansFerreiraPrado2017} is likely driven by funds whose investment restrictions prevent them from acting on the information they receive.

Our paper is most closely related to \cite{Honkanen2025}, who uses hand-collected U.S.\ SEC filings on the ten largest fund families to study securities lending and trading at the quarterly frequency.  \cite{Honkanen2025} finds that active funds reduce their portfolio weights in lent stocks by approximately 1.04 percentage points in the two quarters following a loan, with strong evidence of information spillovers within fund families.  We complement this work along several dimensions.  First, our data cover the universe of German funds rather than a selected set of large families, eliminating sample selection concerns.  Second, we observe lending at the monthly rather than quarterly frequency, providing sharper identification of the timing of information flows.  Third, the European short position disclosure regime gives us an institutional source of variation that has no direct analog in the U.S.\ setting, allowing us to test the information acquisition channel through the interaction of private lending signals and public disclosure events.

Our work also connects to the theoretical literature on information spillovers in securities lending markets.  \cite{Nurisso2026} develops a model in which a securities lender learns about a stock's value by observing whether a short seller borrows shares, and shows that active lenders must raise fees as a commitment device to avoid recalling lent shares.  Our empirical finding that active funds trade on lending information, while passive funds do not, is consistent with his model's prediction that short sellers prefer to borrow from passive lenders who cannot exploit the information.  \cite{ChagueGiovannettiHerskovic2025} identify a distinct but related channel in the Brazilian market: brokers intermediating securities lending transactions leak information about skilled short sellers' directional bets to their other clients.  While their focus is on broker-mediated leakage, our paper demonstrates that the lender itself acquires and trades on the information, even without intermediary involvement.

More broadly, we contribute to the literature on information production and disclosure in financial markets (\citealp{Goldstein_yang:2017, grossman_stiglitz:80}).  The securities lending market represents a setting in which participation itself generates private information, creating a tension between the revenue motive for lending and the trading motive for acting on the information received.  \cite{prado_2015} shows that institutional investors buy shares in response to increases in lending fees, suggesting that expected future lending income is capitalized into prices.  Our paper documents the other side of this trade-off: when the information signal from lending is sufficiently negative, active managers reduce their exposure to avoid capital losses.

We also contribute to the literature on short-sale constraints and price efficiency (\citealp{saffi_sigurdsson:08, prado_saffi_sturgess:14}).  \cite{PaliaSokolinski2024} show that short sellers strategically borrow from passive investors to reduce dynamic short-selling risks, including the risk that active lenders will trade against them.  Our finding that active funds reduce positions in response to lending provides direct evidence for the information leakage concern that motivates short sellers' preference for passive lenders.  Finally, our results relate to the literature on short position disclosure (\citealp{jank_roling_smajlbegovic2020, jones_reed_waller:16}), which shows that mandatory disclosure regimes reduce short interest and price informativeness; we show that the same disclosure events can be used to validate the private information channel from lending to trading.

The organization of the paper is as follows.  Section~\ref{sec:institutional} describes the institutional background and data.  Section~\ref{sec:main_results} reports our main results on position changes by active and passive funds.  Section~\ref{sec:information_channel} tests the information acquisition channel using the disclosure regime.  Section~\ref{sec:family_channels} examines within-family information flows.  Section~\ref{sec:market_structure} analyzes market structure heterogeneity.  Section~\ref{sec:conclusion} concludes.
